\subsection{OpenAI数据集上的精度敏感性}
\label{sec:7-precision}

补证源码检查确认：E1实现虽然输入和存储为FP32，OpenAI旋转与质心GEMM实际采用
TF32乘法、FP32累积，源码注释具有误导性（\cref{sec:4-2}）。主E1前沿保留该历史
实现的测量，另以FP32变体修正cuBLAS开关，重跑37／39个OpenAI配置，batch为1,000，
每配置五轮，在各探针设置间交替变体执行顺序。两者复用同一TF32训练模型，改变的是
add／查询算术，包括数据库编码及路由，而不是独立FP32训练模型的对比。\Cref{tab:7-precision}汇总配对变化。

\par\medskip\noindent\begin{minipage}{\linewidth}
\centering\captionsetup{hypcap=false}
\captionof{table}[匹配配置下的 FP32 敏感性]{FP32敏感性。QPS比值为相同配置下FP32/TF32的几何平均，召回差方向相同。
最后两列在recall 0.97处比较FP32前沿与归档E1的IVF-RaBitQ及CAGRA fp32；
基线没有重新测量。}
\label{tab:7-precision}
\scriptsize\setlength{\tabcolsep}{3pt}
\begin{tabular}{@{}lrrrrr@{}}
\toprule
数据集 & 配对数 & QPS 比值 & $\Delta$Recall & 对 RaBitQ & 对 CAGRA \\
\midrule
OpenAI-1536 & 37 & 0.953 & +0.0003--+0.0005 & 1.086 & 1.679 \\
OpenAI-3072 & 39 & 0.951 & +0.0000--+0.0003 & 1.079 & 1.629 \\
\bottomrule
\end{tabular}

\end{minipage}\par\medskip

相同配置下，FP32吞吐约低5\%，recall高0--0.0005。recall 0.97处，吞吐比同会话
TF32前沿约低3\%，仍为归档RaBitQ的$1.086\times$／$1.079\times$，以及归档
CAGRA fp32的$1.679\times$／$1.629\times$。同会话v69/RaBitQ为1.118／1.114，
历史值为1.130／1.157。这将精度干预与会话波动分开：领先方向仍在、幅度缩小，
跨会话基线比较仍有计时不确定性。两精度变体完整前沿见\cref{fig:7-precision}。

\begin{figure}[tbp]
\centering
\includegraphics[width=\linewidth]{supplement/precision.pdf}
\caption[TF32 与 FP32 精度前沿]{补证精度前沿。实线为同会话JHQ变体，虚线／点线为历史E1基线；
均采用本章前沿及插值规则。}
\label{fig:7-precision}
\end{figure}
